\section{Domain model reference}
\label{sec:models}

This section summarizes the core Django models. Each entry is a one-line purpose statement with links to relevant workflows. Detailed field-level notes appear in later sections.

\subsection{Enum and choice types}
\label{subsec:models:enums}

\begin{longtable}{p{0.30\textwidth}p{0.64\textwidth}}
\toprule
\textbf{Type} & \textbf{Purpose} \\
\midrule
\texttt{ProvisionType} & Normalized provision labels used by students and matching logic (see Section~\ref{sec:ingestion}). \\
\texttt{ExamVenueProvisionType} & Provision capabilities used for venue matching and allocations. \\
\texttt{VenueType} & Categorical labels for venue classification and reporting. \\
\texttt{ExamTypeChoices} & Exam modality labels (on-campus vs online). \\
\texttt{DietChoices} & Predefined diet windows for scheduling (expand as new cycles are added). \\
\texttt{SlotChoices} & Morning vs evening availability slots. \\
\texttt{InvigilatorQualificationChoices} & Qualification labels for invigilator capability tracking. \\
\texttt{InvigilatorRestrictionType} & Restriction labels used when scheduling invigilators. \\
\texttt{AccessibilityFeatures} & Venue accessibility metadata. \\
\bottomrule
\end{longtable}

\subsection{Core domain tables}
\label{subsec:models:core}

\begin{longtable}{p{0.30\textwidth}p{0.64\textwidth}}
\toprule
\textbf{Model} & \textbf{Purpose} \\
\midrule
\texttt{Diet} & Represents an exam diet window with start/end dates and restriction cutoffs. \\
\texttt{Exam} & Stores exam metadata such as course code, school, and student counts. \\
\texttt{Venue} & Represents a physical location with capacity, type, and capabilities. \\
\texttt{ExamVenue} & Links an exam to a venue and captures timing/provision capability context. \\
\texttt{Student} & Stores student identifiers and names. \\
\texttt{StudentExam} & Links students to exams and optionally to matched exam venues. \\
\texttt{Provisions} & Stores student provisions and extra-time notes per exam. \\
\bottomrule
\end{longtable}

\subsection{Notifications and announcements}
\label{subsec:models:notifications}

\begin{longtable}{p{0.30\textwidth}p{0.64\textwidth}}
\toprule
\textbf{Model} & \textbf{Purpose} \\
\midrule
\texttt{Notification} & Tracks system-generated events for admin or invigilator attention. \\
\texttt{Announcement} & Admin-authored messages displayed in the UI. \\
\texttt{UploadLog} & Records upload history for traceability and reporting. \\
\bottomrule
\end{longtable}

\subsection{Invigilator management}
\label{subsec:models:invigilators}

\begin{longtable}{p{0.30\textwidth}p{0.64\textwidth}}
\toprule
\textbf{Model} & \textbf{Purpose} \\
\midrule
\texttt{Invigilator} & Profile for invigilators, linked to auth users when present. \\
\texttt{InvigilatorQualification} & Assigns qualifications to invigilators (unique per invigilator). \\
\texttt{InvigilatorRestriction} & Records diet-specific restrictions per invigilator. \\
\texttt{InvigilatorAvailability} & Tracks availability by date/slot. \\
\texttt{InvigilatorAssignment} & Connects invigilators to exam venues with shift timing and role. \\
\bottomrule
\end{longtable}

\subsection{Where to look next}
\label{subsec:models:next}

\begin{itemize}
  \item Model behaviors and signals: see Section~\ref{sec:signals} (to be added).
  \item API serialization of these models: see Section~\ref{sec:api} (to be added).
  \item Ingestion mapping into these models: see Section~\ref{sec:ingestion}.
\end{itemize}
